\section{Problem setting and inferential scope}
\label{sec:scope}

At a tree node, split selection compares both variables and thresholds. In a
CART-style learner, if feature $j$ has $|C_j|$ admissible thresholds, then the
best observed score for that feature is an extreme statistic over $|C_j|$
opportunities \citep{Breiman1984ClassificationAndRegressionTrees}. Under a
complete null, features with larger candidate sets therefore look more
favorable by chance, which is the core high-cardinality bias documented in the
tree-induction and random-forest importance literature
\citep{LohShih1997SplitSelectionMethods,Shih2004NoteOnSplitSelectionBiasClassificationTrees,KimLoh2001ClassificationTreesUnbiasedMultiwaySplits,Loh2002RegressionTreesUnbiasedVariableSelectionInteractionDetection,StroblBoulesteixAugustin2007UnbiasedSplitSelectionGini,Strobl2007BiasRFVariableImportance}.
Schematically, the greedy rule is
%
\begin{equation}
  (\widehat j,\widehat c) \in \argmax_{j,\ c \in C_j} S_{j,c},
\end{equation}
%
so variables with larger candidate sets receive more chances to attain an
extreme null score.

Conditional inference changes the object being analyzed
\citep{Hothorn2006UnbiasedRecursivePartitioning}. Stage~A asks whether any
candidate feature is associated with the response at the current node; Stage~B
then searches for a split within the selected feature. The underlying
permutation-statistics viewpoint traces to \citet{StrasserWeber1999PermutationStatistics}.
For this paper, the important point is the separation itself: once variable
screening is isolated from threshold optimization, one can formulate a
fixed-node screening problem that is distinct from the full adaptive tree.
The corresponding two-stage schematic is
%
\begin{equation}
  \widehat j \in \argmin_{j \in F_t} p_{t,j},
  \qquad
  \widehat c \in \argmax_{c \in C_{t,\widehat j}} S_{\widehat j,c},
\end{equation}
%
which separates the inferential question ``is any feature associated with the
response at this node?'' from the downstream optimization over thresholds.

That fixed-node screening problem is the inferential target studied here. At a
pre-specified node, Stage~A tests whether any feature in the tested family is
associated with the response under a nodewise complete permutation null. This
is narrower than post-selection inference for the selected feature or
threshold, and narrower than global error control for an end-to-end ranking
workflow \citep{Berk2013ValidPostSelectionInference,Lee2016ExactPostSelectionInferenceLasso,LeebPotscher2015ConfidenceIntervalsAfterModelSelection,BenjaminiHochberg1995FDR,BarberCandes2015Knockoffs,MeinshausenBuhlmann2010StabilitySelection}.
Once Stage~A is embedded in a tree or forest, later nodes and Stage~B become
adaptive because they reuse labels after earlier data-dependent decisions.

Within that fixed-node target, multiplicity is handled conservatively with
Bonferroni across features. The same broad max-statistic idea also appears in
the paper's multi-selector mode, where multiple Stage~A scores are aggregated
within a feature \citep{WestfallYoung1993ResamplingBasedMultipleTesting}. The
paper uses these tools only for the fixed-node decision event; it does not
claim partial-null featurewise inference for the selected variable.

The broader systems context still matters because a usable tree learner needs
more than a calibrated screening rule. Random forests, extra trees, and modern
boosting systems all rely on resampling, subsampling, or approximate search to
stay practical at scale
\citep{Breiman2001RandomForests,Geurts2006ExtraTrees,Friedman2001GBM,ChenGuestrin2016XGBoost,Ke2017LightGBM,Prokhorenkova2018CatBoost}.
Conditional-inference trees face the same pressure. Early stopping, candidate
ordering, threshold subsampling, feature muting, and forest resampling are
therefore part of the practical algorithm, but they should not be conflated
with the fixed-$B$ calibration target.
